\section{E. One-parameter subgroups, envelopes, and completions}

The manuscript's ``metabelian'' class in C4, section 50.4,
is the class of finitely generated virtually abelian groups.
It is not the class of all groups of modern derived length two.
Write $E_L(\Gamma)$ for a finite lattice model,
$\operatorname{Env}(\Gamma)$ for their inverse limit,
$\widehat\Gamma_{\rm const}$ for the affine inverse limit
of constant finite quotients, and $\Gamma_k^{\rm alg}$,
$\Gamma_k^{\rm un}$ for the full proalgebraic and prounipotent
completions. Constant finite schemes remain etale at primes
dividing their orders. In section 51 the manuscript's full
GL group has finite quotient of order 24; the requested
orientation-preserving group $S$ has quotient of order 12.

\begin{statement}{E.1}{Proved}{Characteristic-zero existence and uniqueness}
Let $A$ be a $\mathbb Q$-algebra, $G\hookrightarrow\operatorname{GL}_{r,A}$
closed, and $g\in G(A)$ fibrewise ordinary unipotent.
There is a unique homomorphism $\mathbb G_{a,A}\to G$
taking $1$ to $g$, namely $\exp(T\log g)$.
Uniqueness holds for every separated group scheme over $A$.
Over a characteristic-zero field, existence also holds for
a separated group scheme locally of finite presentation
when the section geometrically lies on an additive subgroup.
Nonlinear existence over an arbitrary $\mathbb Q$-algebra
is not established here.
\end{statement}
\begin{proof}
For $N=g-I$, the entries of $N^r$ lie in the nilradical.
A finitely generated ideal of nilpotent elements is nilpotent:
if the generators have exponents $a_i$, monomials of degree
greater than $\sum(a_i-1)$ vanish. Thus $N$ is nilpotent.
Logarithm and exponential are finite polynomials and satisfy
the group law. Every defining equation of $G$ vanishes at
all nonnegative integral values, since $u(m)=g^m$.
A polynomial of degree $d$ over a $\mathbb Q$-algebra
vanishing at $0,\ldots,d$ is zero by the invertible
Vandermonde determinant. Hence $u$ lands in $G$.

For two homomorphisms into a separated group, the equalizer
is closed in $\mathbb A^1_A$. Its defining polynomials vanish
at all these integers, so the equalizer is the whole line.
Over a field, the two pullbacks of a geometric homomorphism
to $\overline k\otimes_k\overline k$ agree by this uniqueness.
Faithfully flat descent gives existence over $k$.
\end{proof}

\begin{statement}{E.2}{Refuted}{Positive characteristic requires more than ordinary unipotence}
Over a field of characteristic $p$, a matrix $g$ lies at
parameter one on an additive subgroup exactly when
$(g-I)^p=0$; the homomorphism need not be unique.
Even geometric membership in an additive subgroup need
not descend over an imperfect field.
\end{statement}
\begin{proof}
A Jordan block of size $p+1$ has order $p^2$ and cannot
be a value of $\mathbb G_a$, which is killed by $p$.
Likewise in the length-two Witt group, with law

$$
(a,b)+(c,d)=
 \left(a+c,b+d-\sum_{i=1}^{p-1}\binom pi a^ic^{p-i}/p\right),
$$

one has $p(1,0)=(0,1)$, so $(1,0)$ has order $p^2$.
These examples do not satisfy the stronger geometric
hypothesis and are not counterexamples to that hypothesis.

A counterexample which does satisfy it is over
$k=\mathbb F_p(t)$:

$$
U:\ y^p=x-t^{1-p}x^p,\qquad g=(t,0).
$$

The equation is additive and its $x$ derivative is $-1$,
so $U$ is smooth of dimension one.
For $b^p=-t^{1-p}$ the map
$z\mapsto(z^p,z+bz^p)$ is a geometric isomorphism
$\mathbb G_a\to U$, inverse $z=y-bx$.
If polynomials $P,Q$ with zero constants define a map
$\mathbb G_a\to U$, then $Q^p=P-t^{1-p}P^p$.
For $\deg P>0$, leading coefficients would make
$-t^{1-p}$ a $p$th power, impossible by its valuation
$1-p$. Thus $P=Q=0$; no such map takes $1$ to $g$.

For matrices, necessity follows from $g^p=I$.
If $N^p=0$, the logarithm and exponential truncated
below degree $p$ have invertible denominators and give
the homomorphism. Nonuniqueness already follows from
the maps $T$ and $T^p$ on $\mathbb G_a$.
\end{proof}

\begin{statement}{E.3}{Refuted}{A smooth obstruction over every Z/p-squared}
Over $\mathbb Z/p^2$, geometric fibrewise additive
unipotence, even with order $p$, need not give an
integral additive subgroup.
\end{statement}
\begin{proof}
Put $C_p(X,Y)=((X+Y)^p-X^p-Y^p)/p\in\mathbb Z[X,Y]$
and give $\mathbb A^1_{\mathbb Z/p^2}$ the law
$x*y=x+y+pC_p(x,y)$.
The integral cocycle identity for $C_p$ proves associativity;
changes of its arguments by multiples of $p$ disappear
after the outer multiplication by $p$.
The inverse is $-x-pC_p(x,-x)$.
The special fibre is $\mathbb G_a$, and telescoping gives

$$
[p]_G(T)=p(T-T^p).
$$

Thus $1$ has order $p$.
If a homomorphism $P:\mathbb G_a\to G$ takes $1$ to $1$,
write $B=P\bmod p$ and let $c_1$ be its linear coefficient.
Compatibility with multiplication by $p$ yields

$$
\bar c_1T=B-B^p.
$$

But $B(0)=0$, $B(1)=1$, and the right side has degree
$p\deg B>1$, contradiction.
At two the law is $x+y+2xy$ over $\mathbb Z/4$.
The replacement verifies the identities symbolically,
all finite associativity triples at $p=2,3,5$, and
all normalized additive polynomials of degree at most
12: respectively $8,9,5$ candidates and no survivors.
The proof excludes all degrees.
\end{proof}

\begin{statement}{E.4}{Proved}{The divided-power criterion}
For a torsion-free $\mathbb Z$-algebra $A$, a closed linear
group $G/A$, and $g\in G(A)$ with $g-I$ nilpotent,
put $X=\log g$. An additive subgroup taking $1$ to $g$
exists exactly when every $X^j/j!$ is integral.
It is then uniquely $\exp(TX)$.
\end{statement}
\begin{proof}
E.1 forces that polynomial after rationalization,
so each coefficient must be integral if the map exists.
Conversely the coefficients give an integral matrix
polynomial and inverse $\exp(-TX)$.
The group identity and defining equations of $G$ hold
rationally and therefore integrally by torsion-freeness.
The same argument proves uniqueness.
\end{proof}

\begin{statement}{E.5}{Refuted}{The signed implication needs integral X}
For rational nilpotent $X$, integrality of
$X^2/2\pm X^3/6$ does not imply integrality of the summands,
even if $2X$ is integral and primitive.
The same failure holds with $X^3/2$.
The version explicitly assuming integral $X$ is true.
\end{statement}
\begin{proof}
Take

$$
X=\begin{pmatrix}0&3/2&-2&0\\0&0&4&0\\0&0&0&1/2\\0&0&0&0\end{pmatrix}.
$$

Then $A=X^2/2=3E_{13}-E_{14}/2+E_{24}$ and
$B=X^3/6=E_{14}/2$.
Both $A\pm B$ and $A\pm3B$ are integral; neither summand
is integral. However $\exp X$ is not integral, so this
does not refute a hypothesis requiring that too.

If $X$ is integral, away from two average the signed
sums. At two use $A+B=A(I+X/3)$, whose last factor has
an integral inverse by nilpotence.
For the denominator-two alternative use $A(I+X)$.
This proves the corrected assertion.

For C4 scan 13, formula (20), retain the separate readings
$M_5(A)$ and $M_5(K)$.
If the terminal conclusion is merely $X\in M_5(A)$,
the first reading is tautological; the second is false
from higher-power conditions alone, as $X=E_{12}/2$
shows. E.7 addresses the stronger divided-power reading
with all geometric-fibre hypotheses retained.
\end{proof}

\begin{statement}{E.6}{Refuted}{Integral exponential is weaker than divided powers}
Integral $\exp(\pm X)$ does not imply integral divided
powers. The minimum rank for this weaker failure is three.
\end{statement}
\begin{proof}
For $J=E_{12}+E_{23}$, put $X=J-J^2/2$.
Then $\exp X=I+J$ and $\exp(-X)=I-J+J^2$ are integral,
but $X$ and $X^2/2$ are not.
In rank two, $X^2=0$, so integral $\exp X$ forces integral
$X$. The rank-three example has order four modulo two
and fails the manuscript's geometric additive condition.
\end{proof}

\begin{statement}{E.7}{Refuted}{A sharp bound and a genuine rank-five counterexample}
For a torsion-free $\mathbb Z_{(p)}$-algebra, if $N^r=0$,
$N^p\in pM_r(A)$, and $r\le p^2$, then $I+N$ lies on
an additive subgroup. The bound is sharp at rank $p^2+1$.
Over $\mathbb Z$ the minimum rank for failure with every
geometric-fibre hypothesis retained is five.
One example is the shift with weights $1,6,1,6$.
\end{statement}
\begin{proof}
Use

$$
\exp(T\log(I+N))=\sum_{k<r}\binom TkN^k.
$$

For $k=ap+b<p^2$, $0\le b<p$, the numerator $N^k$
is divisible by $p^a$, while $v_p(k!)=a$.
Thus all coefficients are integral.

For sharpness take $p^2$ consecutive shift weights,
equal to $p$ in positions divisible by $p$ and one elsewhere.
Every block of $p$ contains a multiple of $p$.
The product of all weights has valuation $p$, but
$v_p((p^2)!)=p+1$, so the top polynomial coefficient
has valuation $-1$.

For the global example,
$N=E_{12}+6E_{23}+E_{34}+6E_{45}$.
Then $N^2\in2M_5(\mathbb Z)$, $N^3\in3M_5(\mathbb Z)$,
and $N^5=0$. E.2 proves additive membership in every
positive-characteristic fibre, and E.1 in the generic fibre.
But the top coefficient is $N^4_{15}/4!=36/24=3/2$.
Its logarithm is nevertheless integral:

$$
\begin{pmatrix}
0&1&-3&2&-9\\0&0&6&-3&12\\0&0&0&1&-3\\
0&0&0&0&6\\0&0&0&0&0
\end{pmatrix}.
$$

The divided-power integrality flags in degrees zero to
four are true, true, false, true, false.
In rank at most four the geometric hypotheses imply
$N^p\equiv0\pmod p$ for every prime and $N^r=0$.
The local bound applies at every prime; a rational
coefficient integral in every $\mathbb Z_{(p)}$ is integral.
This proves minimality.
\end{proof}

\begin{statement}{E.8}{Proved}{Finite differences and bounded exhaustive checks}
A rational polynomial is integer-valued exactly when
its expansion in $\binom Tj$ has integral coefficients.
The integers $0,1,\ldots$ form a $p$-ordering, with
factorial valuations $v_p(j!)$.
For all positive-length words of lengths at most 12
in weights $\{1,p\}$, the replacement checks all 8190
words for each prime and obtains

$$
\begin{array}{c|rrr}
p&\text{admissible}&\text{negative top valuations}&
\text{first failure degree}\\\hline
2&984&604&4\\3&3734&30&9\\5&7172&0&\text{none}\le12.
\end{array}
$$

Admissible means every consecutive block of $p$ weights
contains $p$. A separate 5292-member signed family
has exactly 46 counterexamples.
\end{statement}
\begin{proof}
Newton's formula is
$f(T)=\sum_j(\Delta^jf)(0)\binom Tj$.
Integral values give integral differences, and the binomial
polynomials are integer-valued also at negative integers.
The product $a(a-1)\cdots(a-j+1)$ is divisible by $j!$,
with equality of valuation at $a=j$.
This proves the special $p$-ordering statement; compare
\cite[Theorem 2, equation (1)]{E_Bhargava2009}.

For a bit word of length $d$ with $a$ weights equal to $p$,
the top valuation is $a-v_p(d!)$. Enumerating every word
gives the displayed counts and the supplied degree tables.
The signed family has
$X_{12}=a/2$, $X_{23}=b$, $X_{34}=c/2$, $X_{13}=d$,
with $1\le a,c\le7$, $1\le b\le12$, $-4\le d\le4$.
The relevant entries are $ab/4,dc/4,bc/4$ and $abc/24$.
Exact divisibility tests give 46 witnesses.
The first ordered by height $\max(a,2b,c,2|d|)$ and
then parameters is $(1,4,3,-2)$, of height eight.
The program independently verifies formal-exponential
and binomial recurrences through degree 12.
These finite families do not exhaust all nilpotent matrices.
\end{proof}

\begin{statement}{E.9}{Proved}{Positive powers over the integers}
For affine finitely presented $G/\mathbb Z$ and
$g\in G(\mathbb Z)$, the following are equivalent:
all geometric fibres are quasi-unipotent; the generic
geometric fibre is quasi-unipotent; some positive power
of $g$ is the value at one of an integral additive subgroup.
No flatness of $G$ is required.
\end{statement}
\begin{proof}
The first implication is immediate.
A generic geometric additive map for a power descends
to $\mathbb Q$ by E.1 and spreads to $\mathbb Z[1/N]$
using its finitely many coordinates and identities.
E.10 extends a parameter rescaling integrally.
Its value at one agrees generically with the corresponding
power of $g$, hence integrally by schematic density
and separatedness. Specialization proves the remaining
implication.

For a matrix $I+N$ with $N^d=0$, take $m=(d-1)!$.
Every $\binom{mT}{j}$ for $j<d$ is integral, since the
factor $mT$ supplies the denominator $j!$.
The binomial sum is a homomorphism with value $g^m$.
Equations of a closed subgroup hold at all integer
parameters and therefore as polynomials over $\mathbb Z$.
\end{proof}

\begin{statement}{E.10}{Proved}{The denominator-clearing lemma}
Let $S$ be quasi-compact, $U=S[1/N]$ schematically dense,
and $G/S$ affine of finite type.
If $f_U:\mathbb A^1_U\to G_U$ has value at zero extending
to a section over $S$, then $f_U(N^bT)$ extends uniquely
for some $b$. A homomorphism extends as a homomorphism.
The parameter is multiplied, not raised to a power.
\end{statement}
\begin{proof}
On a finite affine cover of $S$, choose finite coordinate
generators. Their pullbacks are
$e^*(z)+\sum_{j\ge1}a_jT^j$ with $a_j\in A[1/N]$.
One exponent makes every $N^{bj}a_j$ integral.
The relations hold by the injection $A[T]\to A[1/N][T]$.
The same injection gives uniqueness and gluing;
its two-variable version gives the group identity.
Evaluation at zero explains the necessary constant-term
condition.
\end{proof}

\begin{statement}{E.11}{Open}{The remaining positive-power scope}
The unrestricted non-affine quasi-compact-base assertion
is not proved here. An exact residual is a separated
finitely presented non-affine $G/\mathbb Z$ with
quasi-unipotent generic section.
A proved extension is: for noetherian $A$ and a specified
closed finitely presented $G\hookrightarrow\operatorname{GL}_{r,A}$,
every fibrewise quasi-unipotent section has a positive
power on an additive subgroup.
\end{statement}
\begin{proof}[Partial result and gap]
Choose a common unipotent power at the finitely many
minimal-prime fraction fields. The entries of its
$r$th nilpotent power lie in the nilradical, hence some
higher power vanishes. E.9's binomial construction
gives a GL-valued homomorphism. Rationally it lies in
$G$ by E.1. The finitely many coefficients of the pulled-back
defining equations are integer torsion with zero constants.
Rescaling by a common annihilator kills them and makes
the map land in $G$.

For a non-affine target the finite global coordinates
used in E.10 are unavailable. Scaling does not itself
supply an affine target chart. Abelian schemes do not
provide counterexamples: a characteristic-zero additive
map to an abelian variety is constant, so the relevant
generic power is identity and separatedness extends
that equality. The divided-power examples refute
unpowered statements and do not settle this residual.
\end{proof}

\begin{statement}{E.12}{Proved}{The finite lattice model}
For $1\to L\to\Gamma\to H\to1$, with $L=\mathbb Z^r$
and $H$ finite, there is a smooth affine marked model
$E_L(\Gamma)$ with underlying scheme
$\coprod_H\mathbb A^r$ and exact sequence

$$
1\to L\otimes\mathbb G_a\to E_L(\Gamma)\to H_{\mathbb Z}\to1.
$$

\end{statement}
\begin{proof}
Choose a normalized section $s(h)$ and its cocycle
$c(h,k)=s(h)s(k)s(hk)^{-1}\in L$.
On the disjoint union put

$$
(v,h)(w,k)=(v+hw+c(h,k),hk).
$$

The cocycle identity proves associativity; the inverse is
$(-h^{-1}(v+c(h,h^{-1})),h^{-1})$.
The ring is $\prod_H\operatorname{Sym}_{\mathbb Z}(L^*)$.
The quotient has a scheme section, so is even surjective
on sections, with locally constant component labels understood.
The marking is $ls(h)\mapsto(l,h)$.

For $s'(h)=b(h)s(h)$,
$c'=c+b(h)+hb(k)-b(hk)$ and the marked isomorphism is
$(v,h)\mapsto(v+b(h),h)$.
A homomorphism $f:\Gamma\to G(S)$ together with a
conjugation-equivariant additive map
$u:L\otimes\mathbb G_a\to G$ agreeing on $L$ extends
uniquely by $(v,h)\mapsto u(v)f(s(h))$.
Substitution proves its law and uniqueness.
This is the precise pushout property.
\end{proof}

\begin{statement}{E.13}{Proved}{The integral envelope and its universal range}
For finitely generated virtually abelian $\Gamma$, let
$\mathcal L$ be its normal finite-index free lattices.
Then

$$
\operatorname{Env}(\Gamma)=\varprojlim_{L\in\mathcal L}E_L(\Gamma),
 \qquad
 \mathcal O(\operatorname{Env}(\Gamma))
 =\varinjlim_L\prod_{\Gamma/L}\operatorname{Sym}_{\mathbb Z}(L^*).
$$

For affine finitely presented $G/\mathbb Z$, restriction
is a bijection from envelope morphisms to fibrewise
quasi-unipotent homomorphisms $\Gamma\to G(\mathbb Z)$.
The same holds over a noetherian base for a target with
a specified closed linear embedding.
The finite-presentation Hom formula is a filtered colimit.
\end{statement}
\begin{proof}
Remove torsion from a finite-index abelian subgroup and
intersect its finitely many conjugates to obtain $L_0$.
The lattices $nL_0$ are cofinal.
Inclusions give the transition maps by E.12 and hence
the affine inverse limit.

Over $\mathbb Z$, apply E.9 to a lattice basis and take
a common multiple. The resulting additive maps commute:
their commutators vanish on all integer pairs, a
schematically dense set. The same polynomial argument
gives equivariance under $\Gamma$.
E.12 therefore supplies a finite-stage extension.
Conversely a map to a finitely presented affine target
factors through a finite stage: lift finite generators
and relations, then refine to enforce the Hopf identities.
Every marked element has a positive power in that
stage's lattice. Uniqueness follows from integer-lattice
density on each component.

Over a noetherian base with linear target, use E.11.
Commutation and equivariance hold rationally by E.1;
their finitely many defect coefficients are integer
torsion with zero constants. A common rescaling kills
them. For uniqueness, rational equality leaves finitely
many torsion coefficient differences; a lattice refinement
kills the nonconstant terms, and the constants vanish
on marked elements. Thus the inverse-limit maps agree.
The nonlinear general-base assertion is left to E.27.
\end{proof}

\begin{statement}{E.14}{Proved}{Kernel, flatness, and special fibres}
For virtual rank $r$, the envelope is flat over $\mathbb Z$,
and its kernel over $\widehat\Gamma_{\rm const}$ is

$$
K_r=\operatorname{Spec}\bigl(\mathbb Z+
       (X_1,\ldots,X_r)\mathbb Q[X_1,\ldots,X_r]\bigr).
$$

Its rational fibre is $\mathbb G_a^r$, its characteristic-$p$
fibre is trivial, and
$\operatorname{Env}(\Gamma)_{\mathbb F_p}
=\widehat\Gamma_{{\rm const},\mathbb F_p}$.
For $K_1$, points are compatible sequences $x_n=mx_{mn}$.
They equal $A$ for a $\mathbb Q$-algebra, and zero for
$A=\mathbb Z,\mathbb Z_p$, or a ring killed by an integer.
\end{statement}
\begin{proof}
Finite-stage rings are finite products of polynomial
rings and hence flat; filtered colimits preserve flatness.
Identity-component transitions are $T_n=(m/n)T_m$.
Putting $T_{n,i}=X_i/n$ identifies the kernel ring with
$\bigcup_n\mathbb Z[X_1/n,\ldots,X_r/n]$,
which is the displayed ring: sufficiently divisible $n$
supplies any finite list of rational nonconstant coefficients.
Its positive-degree ideal is divisible by every prime.

For the full special fibre, refinement from $n$ to $pn$
sends an old coordinate to an integer translation plus
$p$ times a new one. Modulo $p$, every old polynomial
becomes constant on refined components. The colimit
is exactly the locally constant function ring on
$\widehat\Gamma$, with injective finite quotient function
maps. It is reduced, zero-dimensional and usually
not of finite type.
The point assertions follow from arbitrary divisibility;
over a $\mathbb Q$-algebra the sequence is uniquely $x_1/n$.
Thus ``characteristic zero'' alone cannot imply $K_1(A)=A$.
\end{proof}

\begin{statement}{E.15}{Refuted}{Fpqc does not become fppf at the inverse limit}
The envelope map to $\widehat\Gamma_{\rm const}$ is
faithfully flat affine, hence an fpqc cover, but need
not be an epimorphism of fppf sheaves.
Failure already occurs for $\Gamma=\mathbb Z$
over $\mathbb Z_{(p)}$, for every prime $p$.
\end{statement}
\begin{proof}
A prime of $C_{\rm lc}(\widehat\Gamma,\mathbb Z)$ determines
a profinite point and a prime of $\mathbb Z$.
The local ring is $\mathbb Q$ or $\mathbb Z_{(p)}$:
locally constant functions become constant on a clopen
neighbourhood of that point.
The envelope algebra localized there is a filtered
colimit of polynomial rings with affine transitions
having nonzero integer linear part. It is torsion-free
over $\mathbb Z_{(p)}$ and thus flat.
E.14 gives nonempty characteristic-$p$ fibres; generically
the transitions are isomorphisms of polynomial algebras.
This proves faithful flatness.

Choose $a\in\widehat{\mathbb Z}$ with
$a_p=\sum_{j\ge1}p^{j!}$ and other components zero.
This $p$-adic number is transcendental.
For if $F(a_p)=0$, its distinct integral partial sums
$u_N$ give $p^{(N+1)!}\mid F(u_N)$, while
$|F(u_N)|\le C N^{\deg F}p^{(\deg F)N!}$.
For large $N$, the nonzero values cannot satisfy both.

A hypothetical fppf lift gives, after localizing a
covering member at a point above $p$, a noetherian local
faithfully flat essentially finite-type algebra $B$
and a coordinate $x$ with $x-a_n\in p^nB$ for all $n$.
Then $dx\in\bigcap_n p^n\Omega_{B/\mathbb Z_{(p)}}=0$
by Krull intersection.
Every minimal prime avoids $p$, since $p$ is a
nonzerodivisor. In each finitely generated characteristic-zero
fraction field the image of $x$ has zero differential,
hence is algebraic over $\mathbb Q$: a transcendental
element belongs to a transcendence basis and has nonzero
differential in characteristic zero.
Multiply the finitely many annihilating integral polynomials
and raise to a power to kill the nilpotent value.
This gives nonzero $F\in\mathbb Z[T]$ with $F(x)=0$ in $B$.
Faithful-flat ideal contraction now gives
$F(a_n)\in p^nB\cap\mathbb Z_{(p)}=p^n\mathbb Z_{(p)}$.
Passing to the limit contradicts transcendence.
The finite lattice quotients of E.12 remain fppf exact.
\end{proof}

\begin{statement}{E.16}{Proved}{Relative and full completions are different}
In characteristic zero, $E_L(\Gamma)_k$ is the relative
unipotent completion for the finite quotient $\Gamma\to H$.
The envelope is their inverse limit, with marked maps

$$
\Gamma_k^{\rm alg}\to\operatorname{Env}(\Gamma)_k
                  \to\Gamma_k^{\rm un}.
$$

For $\mathbb Z$,

$$
\operatorname{Env}(\mathbb Z)_{\mathbb Q}
   =\mathbb G_a\times\widehat{\mathbb Z}_{\rm const},
 \quad n\mapsto(n,\widehat n).
$$

Its ordinary unipotent completion is $\mathbb G_a$.
Relative completion for $n\mapsto2^n$ into $\mathbb G_m$
is $\mathbb G_a\times\mathbb G_m$, marked by $(n,2^n)$.
There is no marked morphism in either direction between
the last group and the envelope. This refutes Lead 3.
\end{statement}
\begin{proof}
Use the relative universal problem of
\cite[Section 1; Proposition 11.1]{E_Hain1998}.
In an extension of $H$ by a prounipotent group, the
lattice images have commuting logarithms in every
finite-dimensional quotient. They extend uniquely
linearly and then by exponential. Uniqueness gives
equivariance, and E.12 gives the unique extension on
all components. This proves the claimed universal property.
The full representation completion maps to every finite
envelope stage; the same lattice argument maps the
envelope to unipotent completion.

For $\mathbb Z$, the coordinate $X=a+nT$ on component
$a\bmod n$ identifies the rational stage with
$\mathbb G_a\times(\mathbb Z/n)$, compatibly in $n$.
A character from the envelope to $\mathbb G_m$ factors
through a finite stage, is trivial on $\mathbb G_a$,
and has finite marked image. Hence $2^n$ does not extend.

For the relative infinite quotient, in an algebraic
extension of $\mathbb G_m$ by a unipotent group write
a lift of $2$ as $g_sg_u$.
The diagonalizable cyclic closure of $g_s$ maps onto
$\mathbb G_m$ and has kernel both diagonalizable and
unipotent, hence trivial. It is therefore $\mathbb G_m$.
It centralizes $g_u$, so
$(t,z)\mapsto\exp(t\log g_u)g_s(z)$ supplies the desired
map. The points $(n,2^n)$ are Zariski dense: a relation
$\sum_jP_j(n)2^{jn}=0$ is impossible by the largest
exponential term. Uniqueness and passage to inverse
limits prove the relative universal property.
An envelope map to this group would extend $2^n$.
A reverse marked map would map a connected group
nontrivially to the constant quotient $C_2$, impossible.
\end{proof}

\begin{statement}{E.17}{Proved}{Integer-valued, divided-power, and Hall models}
For

$$
\operatorname{Int}(\mathbb Z)=
 \{f\in\mathbb Q[T]:f(\mathbb Z)\subset\mathbb Z\},\qquad
 D=\bigoplus_{j\ge0}\mathbb Z\,T^j/j!,
$$

there are additive Hopf maps

$$
\operatorname{Env}(\mathbb Z)\to
 \operatorname{Spec}\operatorname{Int}(\mathbb Z),\qquad
 K_1\to\operatorname{Spec}D.
$$

Moreover $\operatorname{Int}(\mathbb Z)/p
=C_{\rm lc}(\mathbb Z_p,\mathbb F_p)$, while $D/p$
is nonreduced. The rank-one integral Hall model and
the abelian BCH model are additive groups, not the envelope.
\end{statement}
\begin{proof}
The binomial basis and
$\binom{T+U}{j}=\sum_i\binom Ti\binom U{j-i}$
give the first Hopf algebra.
Choose $n$ divisible by ordinary coefficient denominators;
then $f(a+nT)$ is integral on every stage component.
This embeds its ring in the envelope ring and respects
Hopf operations by equality on integer points.
The inclusion $D\subset\mathbb Z+T\mathbb Q[T]$ gives
the second map.

Integer-valued polynomials extend continuously to
$\mathbb Z_p$. Reduction has kernel $p\operatorname{Int}(\mathbb Z)$.
The functions $\binom a{p^i}\bmod p$ are the base-$p$
digits, by expansion of $(1+X)^a$.
Digit functions and finite-field interpolation give
all functions on $\mathbb Z/p^m$, hence every locally
constant function. The corresponding special-fibre
map is $\widehat{\mathbb Z}\to\mathbb Z_p$.
In $D/p$, the nonzero basis vector $\gamma_1=T$ satisfies
$\gamma_1^p=p!\gamma_p=0$.

For a generator $b$ of $\mathbb Z$, collection is
$b^ab^c=b^{a+c}$, so the Hall polynomial is $a+c$.
All brackets vanish in the abelian $\mathbb Z_p$ Lie
algebra, giving the same BCH law. These direct rank-one
calculations require no general Lazard hypothesis.
The envelope maps to its additive finite stage, but
has the additional profinite generic component factor
and the constant special fibres of E.14.
\end{proof}

\begin{statement}{E.18}{Proved}{The marked orientation-preserving modular quotient}
Put

$$
\lambda_0=\sigma\rho\sigma^{-1}\rho=
 \begin{pmatrix}-1&2\\0&-1\end{pmatrix},\qquad
 \lambda_1=\rho\lambda_0\rho^{-1}=
 \begin{pmatrix}-1&0\\-2&-1\end{pmatrix}.
$$

Then $\Pi_0=\langle\lambda_0,\lambda_1\rangle$ is free
of rank two, normal of indices 12 and 24 in SL and GL.
For $S=\operatorname{SL}_2(\mathbb Z)/[\Pi_0,\Pi_0]$,

$$
S=\mathbb Z^2\rtimes H,\qquad H=\operatorname{Dic}_3,
$$

where, in this presentation, $r^6=1$, $s^2=r^3$,
$srs^{-1}=r^{-1}$, and the action is

$$
R=\begin{pmatrix}0&-1\\1&-1\end{pmatrix},\quad
 J=\begin{pmatrix}0&1\\1&0\end{pmatrix}.
$$

The marking is $\rho\mapsto(-e_0,r)$, $\sigma\mapsto(0,s)$,
$\lambda_i\mapsto(e_i,1)$. This group is virtually
abelian but not modern metabelian.
\end{statement}
\begin{proof}
C.2 identifies the projective level-two kernel as free
on the corresponding elementary matrices.
The signed lifts $\lambda_i$ give a complement to
$\{\pm I\}$ in $\Gamma(2)$, hence index 12.
Modulo four their subgroup is

$$
K=\left\{I,
 \begin{pmatrix}3&2\\0&3\end{pmatrix},
 \begin{pmatrix}3&0\\2&3\end{pmatrix},
 \begin{pmatrix}1&2\\2&1\end{pmatrix}\right\}.
$$

Conjugation by $\rho$ sends its generators $a,b$ to
$b,ab$, and $\sigma$ exchanges them.
The kernel modulo two in SL modulo four has order eight
and is generated by $a,b,-I$; the total order is 48.
The inverse image of $K$ thus has index 12 and equals
$\Pi_0$. Finally $\operatorname{diag}(-1,1)$ inverts
both $\lambda_i$, proving GL normality and index 24.

The quotient has the displayed dicyclic presentation
and exactly 12 elements by the index.
The matrices satisfy $R^3=J^2=I$, $JRJ=R^{-1}$.
The marked $\rho,\sigma$ satisfy their modular relations,
and direct multiplication sends the two free basis
elements to the lattice basis. The resulting surjection
has kernel exactly $[\Pi_0,\Pi_0]$.
In the derived group, $(1,-1)$ and $r^2$ have commutator
$(3,0)\ne0$, so the group is not metabelian.
The finite quotient at two independently has derived
orders $48,12,4,1$.
\end{proof}

\begin{statement}{E.19}{Proved}{Coordinates, points, and fibres of S}
The finite model is
$\underline S=\mathbb G_a^2\rtimes H$ with ring
$\prod_H\mathbb Z[X,Y]$: twelve smooth affine planes
in every characteristic.
The envelope ring is

$$
\varinjlim_n
 \prod_{(a,h)\in(\mathbb Z/n)^2\rtimes H}\mathbb Z[T_1,T_2].
$$

For $m=nq$ and $b=a+nc$, transitions are $T_n=c+qT_m$.
Furthermore

$$
\widehat S=\widehat{\mathbb Z}^{\,2}\rtimes H,\quad
 \operatorname{Env}(S)(\mathbb Z)=S,\quad
 \operatorname{Env}(S)(\mathbb Z_p)=\widehat S,\quad
 \operatorname{Env}(S)(\widehat{\mathbb Z})=\prod_p\widehat S.
$$

For the finite model,
$\underline S(\widehat{\mathbb Z})=\prod_p(\mathbb Z_p^2\rtimes H)$;
its common-component subset is $\widehat S$.
The full special fibre is constant $\widehat S$, not
twelve affine planes.
\end{statement}
\begin{proof}
Represent the point in component $(a,h)$ as $(a+nT,h)$.
Its semidirect multiplication gives the Hopf operations
with integral carries and all transition formulas.
These finite quotients are cofinal, proving the completion.
At integral points every stage is the same group $S$,
with identity transition.
At $\mathbb Z_p$ a compatible component label is
$(a,h)\in\widehat S$, and the first lattice coordinate
is forced by $x-a_n\in n\mathbb Z_p^2$ to be $a_p$.
Conversely $(a_p-a_n)/n$ supplies all higher coordinates.
Affine functors preserve products of rings.
E.14 gives the full special fibre.
The marked generators generate quotients of orders
$12n^2$ for $n=1,2,3,4,8,9$, as actually checked.
\end{proof}

\begin{statement}{E.20}{Proved}{The specified relative quotient}
The quotient of relative unipotent completion of
$\operatorname{SL}_2(\mathbb Z)\to H$ by the closed
commutator of its prounipotent kernel is
$\underline S_{\mathbb Q}$, with E.18's marking.
It is not the full integral envelope.
\end{statement}
\begin{proof}
A lift into an extension of $H$ by a vector group kills
$[\Pi_0,\Pi_0]$, hence factors through $S$.
E.16 extends it uniquely to $\underline S_{\mathbb Q}$.
That group itself is such an extension, proving the
universal property. Abelianizing a prounipotent kernel
represents precisely this problem.
The full GL variant has 24 components; the envelope
also retains every lattice refinement.
\end{proof}

\begin{statement}{E.21}{Open}{The full normalizer action remains unfinished}
A sign-consistent construction of all actions on C4
scans 36--41 on one marked integral relative-completion
model is not completed here.
Proved: $\underline S_{\mathbb Q}$ has dimension two
and cannot be $M_{\mathbb Q}$ of dimension six.
B.23 computes the finer three-dimensional parameter
scheme. Inner conjugation acts on $\underline S$,
with kernel the constant $\langle r^3\rangle$
over $\mathbb Z[1/6]$.
\end{statement}
\begin{proof}[Partial assertions and gap]
A centralizing $(v,h)$ must have $h$ in the kernel
$\langle r^3\rangle$ of the lattice action.
Then $v$ is fixed by $R,J$; $\det(R-I)=3$ forces
$v=0$ on the stated base.
This proves the inner action and kernel.
Dimension excludes the proposed group-scheme identification.
The finite checks in E.19 and B.23 do not identify an
external parameter group with a full automorphism
functor or verify the entire sign-bearing action.
The arithmetic quotient in E.28 is also outside its
proved scope.
\end{proof}

\begin{statement}{E.22}{Proved}{An explicit Stallings-type test}
For $F=\langle x,y\rangle$, let
$f(x)=x$, $f(y)=y[x,y]$.
It induces automorphisms of every lower-central quotient,
every pro-$\ell$ completion, and the rational Malcev
completion. It is not onto $F$ or its full profinite
completion. The representation $x\mapsto(12)$,
$y\mapsto(13)$ has image six and, after $f$, image two.
\end{statement}
\begin{proof}
The map is identity on abelianization.
Each graded lower-central piece is generated by
commutators in the generators, so changing a generator
by a weight-two term changes no associated graded term.
Induction on the central series gives an automorphism
of each lower-central quotient.

On the free pro-$\ell$ group it is identity modulo
$\Phi(P)=\overline{P^\ell[P,P]}$.
Its compact image is closed. A proper image would
remain proper in a finite $\ell$-quotient and lie
in a maximal subgroup of index $\ell$, contradicting
Frattini surjectivity.
A finitely generated profinite group is Hopfian:
there are finitely many open normal subgroups of any
given index, and inverse image under a surjective
endomorphism permutes them. Some iterate preserves each,
forcing the kernel into every open normal subgroup.
Thus the map is an automorphism.

On the completed free Lie algebra it sends each generator
to itself plus higher degree. A continuous inverse is
constructed recursively degree by degree.
Finally direct multiplication of the two transpositions
gives $y[x,y]=x$. The image drops from six to two,
which is impossible for a surjective endomorphism.
The finite-group program checks this.
\end{proof}

\begin{statement}{E.23}{Refuted}{The two S3-cover claims differ}
The representation of $\Pi_0$ sending its free basis
to $(12),(13)$ cannot factor through $S$ or its marked
envelope. The modern metabelian quotient
$\Pi_0/\Pi_0''$ retains it but is not virtually abelian.
In contrast, ordinary reduction
$\operatorname{SL}_2(\mathbb Z)\to S_3$ factors through
$S$ and its envelope.
\end{statement}
\begin{proof}
The first images do not commute, while $S$ kills
$[\Pi_0,\Pi_0]$. Since $S_3''=1$, the modern metabelian
quotient retains them.
For a free rank-two group, the universal abelian
covering graph has boundary

$$
\mathbb Z[u^{\pm1},v^{\pm1}]^2
 \xrightarrow{(u-1,v-1)}\mathbb Z[u^{\pm1},v^{\pm1}].
$$

Its kernel is freely generated over the Laurent ring
by $(1-v,u-1)$: divisibility by the prime $u-1$
proves this directly. There are no two-cells.
Thus $F'/F''$ is the additive Laurent polynomial group,
not finitely generated. Every subgroup of a finitely
generated virtually abelian group is finitely generated,
so $F/F''$ is not virtually abelian.

Reduction modulo two kills $\Pi_0$ and factors through
$H/\langle r^3\rangle=S_3$, a retained finite quotient.
Its restriction to $\Pi_0$ is trivial and is therefore
a different representation.
In E.22 the transitive three-letter action pulls back
to a two-element orbit and a fixed point, giving a
disconnected pulled-back degree-three cover.
\end{proof}

\begin{statement}{E.24}{Known}{The two rational theories}
Fibrewise rationalization preserves the discrete
fundamental group of $BS=K(S,1)$ and changes no higher
homotopy groups, so $BS$ is already fibrewise rational.
Its rational local-system derived category is $D(\mathbb Q[S])$.
The fundamental affine group of its schematization is
$S_{\mathbb Q}^{\rm alg}$, whereas

$$
\operatorname{Env}(S)_{\mathbb Q}
 =(\mathbb G_a^2\times\widehat{\mathbb Z}_{\rm const}^{\,2})
       \rtimes H.
$$

All retain the reduction quotient to $S_3$.
The marked maps of $S$ into its envelope and full
proalgebraic completion are injective.
No assertion about the entire schematization being a
classifying stack is made without a goodness hypothesis.
\end{statement}
\begin{proof}[Exact reductions and explicit calculation]
The fundamental-group-preserving rational equivalences
are those of \cite[Theorem 3]{E_RiveraWierstraZeinalian2021}.
There are no higher groups to rationalize in $BS$.
Descent from its contractible universal cover identifies
complexes of local systems with derived $\mathbb Q[S]$-modules.
The schematic fundamental-group statement is
\cite[Corollary 3.3.7]{E_Toen2006}.

At a rational envelope stage, use $X=a+nT$.
The multiplication becomes
$(X,a,h)(Y,b,k)=(X+hY,a+hb,hk)$.
Taking the inverse limit gives the displayed formula
and marking $(v,h)\mapsto((v,\widehat v),h)$.
The finite quotients $(\mathbb Z/n)^2\rtimes H$ separate
all elements; their permutation representations also
separate elements in full proalgebraic completion.
The finite quotient $S\to H\to S_3$ survives throughout.
\end{proof}

\begin{statement}{E.25}{Proved}{Representations separating the theories}
A twelve-dimensional rational representation of $S$
extends to full proalgebraic completion but not its
envelope. The regular local system $\mathbb Q[S]$
is not the restriction of any rational comodule of
an affine group scheme along a marked map from $S$.
\end{statement}
\begin{proof}
Induce $\chi(a,b)=2^a$ from $L=\mathbb Z^2$ to $S=L\rtimes H$.
On basis $e_h$, let $v e_h=\chi(h^{-1}v)e_h$ and let
$H$ permute the basis. The lattice generator has
eigenvalue two. An envelope representation factors
through a finite stage, where a positive power of
each marked element lies on an additive subgroup;
its eigenvalues must be roots of unity.
Thus this representation does not extend to the envelope,
but extends by the definition of full proalgebraic completion.

Every vector in a rational comodule lies in a finite-dimensional
subcomodule. Indeed write
$\delta(v)=\sum_i v_i\otimes a_i$ with independent $a_i$.
The counit puts $v$ in $W=\operatorname{span}(v_i)$.
Projecting the first tensor factor of coassociativity
to $V/W$ gives $\delta(v_i)\in W\otimes B$.
But the orbit of $1$ in $\mathbb Q[S]$ contains the
whole infinite basis $S$, so cannot lie in such a
finite-dimensional invariant space.
\end{proof}

\begin{statement}{E.26}{Proved}{The nilpotent completions of S}
For $L_0=(R-I)\mathbb Z^2$,

$$
S'=L_0\rtimes\langle r^2\rangle,\qquad
 S^{ab}=C_{12},\qquad \gamma_nS=S'\ (n\ge2).
$$

Its lower-central and finite pronilpotent completions
are $C_{12}$; its pro-two and pro-three completions
are $C_4,C_3$; other pro-$\ell$ completions and its
characteristic-zero Malcev completion are trivial.
\end{statement}
\begin{proof}
The columns of $J-I$ lie in $(R-I)\mathbb Z^2$,
whose index is three.
Also $H'=\langle r^2\rangle=C_3$ and $H^{ab}=C_4$.
This computes $S'$ and $S^{ab}=C_3\times C_4$.
Since $R^2-I=(R-I)(R+I)$ and $\det(R+I)=1$,
$[L,\langle r^2\rangle]=L_0$.
Also $[s,r^2]=r^2$, so $[S,S']=S'$.
Every nilpotent quotient therefore factors through
$C_{12}$. Taking Sylow quotients gives the pro-$\ell$
claims. A characteristic-zero unipotent group is nilpotent
and torsion-free, by logarithm, so all such images vanish.
The finite quotients are envelope quotients.
The checked mod-two quotient has lower-central orders
$48,12,12,12$ and derived orders $48,12,4,1$.
\end{proof}

\begin{statement}{E.27}{Open}{Two exact general-base residuals}
Unproved here are:
existence for a nonlinear separated locally finitely
presented group over an arbitrary $\mathbb Q$-algebra
with a geometrically additive section;
and the arbitrary-noetherian-base envelope universal
property for affine targets without a supplied linear
embedding. Specific unresolved instances in this work
are a nonlinear group over $\mathbb Q[\epsilon]/(\epsilon^2)$
and such an affine group over $\mathbb Z[t]$.
These are not claims of literature-wide openness.
\end{statement}
\begin{proof}[Proved range and missing inputs]
E.1 proves all closed linear cases over $\mathbb Q$-algebras
and general groups over characteristic-zero fields.
Nilpotents alone are not an obstruction in that range:

$$
g=\begin{pmatrix}1+\epsilon&1\\0&1\end{pmatrix},\qquad
 u(T)=\begin{pmatrix}
 1+\epsilon T&T+\epsilon(T^2-T)/2\\0&1
 \end{pmatrix}
$$

satisfies $u(T+U)=u(T)u(U)$ and $u(1)=g$.
The program checks these identities modulo $\epsilon^2$.
Without a linear embedding, the missing input is
existence of the interpolating morphism into the
nonlinear target; fieldwise descent alone does not
construct a map on a thickening.
For the general noetherian affine case the missing
input is an appropriate positive-power theorem.
Once additive lattice maps exist, E.13 handles their
finite torsion defects and uniqueness after refinement.
\end{proof}

\begin{statement}{E.28}{Proved}{Discrete Fermat relative completions}
Let $F=\pi_1(\mathbb P^1(\mathbb C)\setminus\{0,1,\infty\})
=\langle x,y\rangle$ and
$q_n:F\to A_n=(\mathbb Z/n)^2$ send meridians to its basis.
Relative completion over $\mathbb Q$ has free prounipotent
kernel of rank $n^2+1$.
Its abelian-kernel quotient has dimension $n^2+1$,
equal to five at $n=2$; its full completion is not
of finite type. It cannot be $M_{\mathbb Q}$.
For $n\ge2$ its abelian-kernel quotient cannot be
$\underline S_{\mathbb Q}$.
The arithmetic/profinite comparison relative to
$S_3\times\widehat{\mathbb Z}^{\times}$ remains unresolved.
\end{statement}
\begin{proof}
The covering graph of the two-petal rose has $n^2$
vertices and $2n^2$ positive edges.
Contracting a spanning tree leaves $n^2+1$ edges,
so $K_n=\ker q_n$ is free of that rank.
This is the connected cover adjoining $n$th roots
of $z$ and $1-z$.
The executed graph check covers $1\le n\le12$.
At two it gives the basis

$$
x^2,\quad yxy^{-1}x^{-1},\quad y^2,\quad
 xyxy^{-1},\quad xy^2x^{-1}.
$$


Let $U_n$ be exponential of the completed free Lie
algebra on that basis.
Assigning logarithms of the generators uniquely
extends a map from $K_n$ to any finite-dimensional
unipotent group, and then to inverse limits.
Thus $U_n$ is its unipotent completion.
Choose representatives $s(a)$ in $F$.
Their conjugations extend to $C_a$ on $U_n$, and put
$c(a,b)=s(a)s(b)s(a+b)^{-1}$ in $U_n$.
They satisfy

$$
C_aC_b=\operatorname{Inn}(c(a,b))C_{a+b},\qquad
 c(a,b)c(a+b,d)=C_a(c(b,d))c(a,b+d).
$$

Hence

$$
(u,a)(v,b)=(uC_a(v)c(a,b),a+b)
$$

defines a proalgebraic extension with the given marking.
Any competing relative extension receives the unique
map from $U_n$; uniqueness forces compatibility with
the conjugations and cocycle, extending it on each
component. This is the relative universal property.

For $s'(a)=b_as(a)$,

$$
C'_a=\operatorname{Inn}(b_a)C_a,\qquad
 c'(a,b)=b_aC_a(b_b)c(a,b)b_{a+b}^{-1}.
$$

The isomorphism to old coordinates is $(u,a)\mapsto(ub_a,a)$,
as direct substitution shows.
Thus there is no dependence on the section.

The free Lie abelianization has dimension $n^2+1$.
For at least two generators, $\operatorname{ad}(X)^mY$
has associative word $X^mY$ with coefficient one,
so nonzero degrees occur arbitrarily far out.
The full kernel is not of finite type.
No square $n^2$ equals five, so the abelian-kernel
dimension is never six, and exceeds two for $n\ge2$.
These are discrete complex comparisons; adjoining
the requested arithmetic quotient requires its own
marking, topology and representation category.
\end{proof}
